% Options for packages loaded elsewhere
\PassOptionsToPackage{unicode}{hyperref}
\PassOptionsToPackage{hyphens}{url}
%
\documentclass[
  ignorenonframetext,
]{beamer}
\usepackage{pgfpages}
\setbeamertemplate{caption}[numbered]
\setbeamertemplate{caption label separator}{: }
\setbeamercolor{caption name}{fg=normal text.fg}
\beamertemplatenavigationsymbolsempty
% Prevent slide breaks in the middle of a paragraph
\widowpenalties 1 10000
\raggedbottom
\setbeamertemplate{part page}{
  \centering
  \begin{beamercolorbox}[sep=16pt,center]{part title}
    \usebeamerfont{part title}\insertpart\par
  \end{beamercolorbox}
}
\setbeamertemplate{section page}{
  \centering
  \begin{beamercolorbox}[sep=12pt,center]{part title}
    \usebeamerfont{section title}\insertsection\par
  \end{beamercolorbox}
}
\setbeamertemplate{subsection page}{
  \centering
  \begin{beamercolorbox}[sep=8pt,center]{part title}
    \usebeamerfont{subsection title}\insertsubsection\par
  \end{beamercolorbox}
}
\AtBeginPart{
  \frame{\partpage}
}
\AtBeginSection{
  \ifbibliography
  \else
    \frame{\sectionpage}
  \fi
}
\AtBeginSubsection{
  \frame{\subsectionpage}
}
\usepackage{amsmath,amssymb}
\usepackage{lmodern}
\usepackage{iftex}
\ifPDFTeX
  \usepackage[T1]{fontenc}
  \usepackage[utf8]{inputenc}
  \usepackage{textcomp} % provide euro and other symbols
\else % if luatex or xetex
  \usepackage{unicode-math}
  \defaultfontfeatures{Scale=MatchLowercase}
  \defaultfontfeatures[\rmfamily]{Ligatures=TeX,Scale=1}
\fi
\usetheme[]{Montpellier}
% Use upquote if available, for straight quotes in verbatim environments
\IfFileExists{upquote.sty}{\usepackage{upquote}}{}
\IfFileExists{microtype.sty}{% use microtype if available
  \usepackage[]{microtype}
  \UseMicrotypeSet[protrusion]{basicmath} % disable protrusion for tt fonts
}{}
\makeatletter
\@ifundefined{KOMAClassName}{% if non-KOMA class
  \IfFileExists{parskip.sty}{%
    \usepackage{parskip}
  }{% else
    \setlength{\parindent}{0pt}
    \setlength{\parskip}{6pt plus 2pt minus 1pt}}
}{% if KOMA class
  \KOMAoptions{parskip=half}}
\makeatother
\usepackage{xcolor}
\IfFileExists{xurl.sty}{\usepackage{xurl}}{} % add URL line breaks if available
\IfFileExists{bookmark.sty}{\usepackage{bookmark}}{\usepackage{hyperref}}
\hypersetup{
  hidelinks,
  pdfcreator={LaTeX via pandoc}}
\urlstyle{same} % disable monospaced font for URLs
\newif\ifbibliography
\usepackage{color}
\usepackage{fancyvrb}
\newcommand{\VerbBar}{|}
\newcommand{\VERB}{\Verb[commandchars=\\\{\}]}
\DefineVerbatimEnvironment{Highlighting}{Verbatim}{commandchars=\\\{\}}
% Add ',fontsize=\small' for more characters per line
\usepackage{framed}
\definecolor{shadecolor}{RGB}{248,248,248}
\newenvironment{Shaded}{\begin{snugshade}}{\end{snugshade}}
\newcommand{\AlertTok}[1]{\textcolor[rgb]{0.94,0.16,0.16}{#1}}
\newcommand{\AnnotationTok}[1]{\textcolor[rgb]{0.56,0.35,0.01}{\textbf{\textit{#1}}}}
\newcommand{\AttributeTok}[1]{\textcolor[rgb]{0.77,0.63,0.00}{#1}}
\newcommand{\BaseNTok}[1]{\textcolor[rgb]{0.00,0.00,0.81}{#1}}
\newcommand{\BuiltInTok}[1]{#1}
\newcommand{\CharTok}[1]{\textcolor[rgb]{0.31,0.60,0.02}{#1}}
\newcommand{\CommentTok}[1]{\textcolor[rgb]{0.56,0.35,0.01}{\textit{#1}}}
\newcommand{\CommentVarTok}[1]{\textcolor[rgb]{0.56,0.35,0.01}{\textbf{\textit{#1}}}}
\newcommand{\ConstantTok}[1]{\textcolor[rgb]{0.00,0.00,0.00}{#1}}
\newcommand{\ControlFlowTok}[1]{\textcolor[rgb]{0.13,0.29,0.53}{\textbf{#1}}}
\newcommand{\DataTypeTok}[1]{\textcolor[rgb]{0.13,0.29,0.53}{#1}}
\newcommand{\DecValTok}[1]{\textcolor[rgb]{0.00,0.00,0.81}{#1}}
\newcommand{\DocumentationTok}[1]{\textcolor[rgb]{0.56,0.35,0.01}{\textbf{\textit{#1}}}}
\newcommand{\ErrorTok}[1]{\textcolor[rgb]{0.64,0.00,0.00}{\textbf{#1}}}
\newcommand{\ExtensionTok}[1]{#1}
\newcommand{\FloatTok}[1]{\textcolor[rgb]{0.00,0.00,0.81}{#1}}
\newcommand{\FunctionTok}[1]{\textcolor[rgb]{0.00,0.00,0.00}{#1}}
\newcommand{\ImportTok}[1]{#1}
\newcommand{\InformationTok}[1]{\textcolor[rgb]{0.56,0.35,0.01}{\textbf{\textit{#1}}}}
\newcommand{\KeywordTok}[1]{\textcolor[rgb]{0.13,0.29,0.53}{\textbf{#1}}}
\newcommand{\NormalTok}[1]{#1}
\newcommand{\OperatorTok}[1]{\textcolor[rgb]{0.81,0.36,0.00}{\textbf{#1}}}
\newcommand{\OtherTok}[1]{\textcolor[rgb]{0.56,0.35,0.01}{#1}}
\newcommand{\PreprocessorTok}[1]{\textcolor[rgb]{0.56,0.35,0.01}{\textit{#1}}}
\newcommand{\RegionMarkerTok}[1]{#1}
\newcommand{\SpecialCharTok}[1]{\textcolor[rgb]{0.00,0.00,0.00}{#1}}
\newcommand{\SpecialStringTok}[1]{\textcolor[rgb]{0.31,0.60,0.02}{#1}}
\newcommand{\StringTok}[1]{\textcolor[rgb]{0.31,0.60,0.02}{#1}}
\newcommand{\VariableTok}[1]{\textcolor[rgb]{0.00,0.00,0.00}{#1}}
\newcommand{\VerbatimStringTok}[1]{\textcolor[rgb]{0.31,0.60,0.02}{#1}}
\newcommand{\WarningTok}[1]{\textcolor[rgb]{0.56,0.35,0.01}{\textbf{\textit{#1}}}}
\setlength{\emergencystretch}{3em} % prevent overfull lines
\providecommand{\tightlist}{%
  \setlength{\itemsep}{0pt}\setlength{\parskip}{0pt}}
\setcounter{secnumdepth}{-\maxdimen} % remove section numbering
%\documentclass[unicode,10pt]{beamer}% 'unicode'が必要
\usepackage{luatexja}% 日本語を使う場合は必須
%\usepackage[japanese]{babel}	
%\usefonttheme[onlymath]{serif}
%\usepackage[T1]{fontenc} %これを使うと、ギリシャ文字が出ない
%\usepackage{textcomp}
%\usepackage[scale=1.0]{tgheros} % Sans serif font
%\usepackage[scaled]{beramono} % Monospace font
%\usepackage{luatexja-otf}
%\renewcommand{\kanjifamilydefault}{\gtdefault}
%\usepackage[haranoaji]{luatexja-preset}

% スライド番号の表示 %%%%%%%%%%%%%%%%%%%
\makeatletter
\setbeamertemplate{footline}{%
  \leavevmode%
  \hbox{%
  \begin{beamercolorbox}[wd=.5\paperwidth,ht=2.25ex,dp=1ex,right]{date in head/foot}%
    \insertframenumber{} \hspace*{2ex} % new version without total frames
  \end{beamercolorbox}}%
  \vskip0pt%
}
\makeatother
% スライド番号の表示 %%%%%%%%%%%%%%%%%%%
\ifLuaTeX
  \usepackage{selnolig}  % disable illegal ligatures
\fi

\title{第2章: 無作為化実験による因果効果の推定}
\subtitle{Elena Llaudet \& Kosuke Imai.\\
Data Analysis for Social Science: A Friendly and Practical
Introduction.}
\author{}
\date{\vspace{-2.5em}2026-03-09}

\begin{document}
\frame{\titlepage}

\hypertarget{starux30d7ux30edux30b8ux30a7ux30afux30c8}{%
\section{2.1
STARプロジェクト}\label{starux30d7ux30edux30b8ux30a7ux30afux30c8}}

\begin{frame}{STARプロジェクトの概要}
\protect\hypertarget{starux30d7ux30edux30b8ux30a7ux30afux30c8ux306eux6982ux8981}{}
\begin{itemize}
\tightlist
\item
  1980年代にテネシー州で実施された「生徒と教師の学習達成率プロジェクト」。
\item
  目的: 学級の大きさが生徒の成績に及ぼす影響を検証すること。
\item
  幼稚園児を小学3年生の終わりまで、以下のいずれかに\textbf{無作為に割り当て}た：

  \begin{itemize}
  \tightlist
  \item
    \textbf{少人数学級} (13人〜17人)
  \item
    \textbf{標準規模学級} (22人〜25人)
  \end{itemize}
\end{itemize}
\end{frame}

\hypertarget{ux51e6ux7f6eux5909ux6570ux3068ux7d50ux679cux5909ux6570}{%
\section{2.2
処置変数と結果変数}\label{ux51e6ux7f6eux5909ux6570ux3068ux7d50ux679cux5909ux6570}}

\begin{frame}{因果関係の方向性}
\protect\hypertarget{ux56e0ux679cux95a2ux4fc2ux306eux65b9ux5411ux6027}{}
\begin{itemize}
\tightlist
\item
  \textbf{因果関係}: 2つの変数の間の原因と結果の関係。
\item
  \textbf{処置変数 (\(X\))}: 変化が発生する変数（原因）。
\item
  \textbf{結果変数 (\(Y\))}:
  処置変数の変化に応答して変化しうる変数（結果）。
\end{itemize}

\[ X \rightarrow Y \]

\begin{itemize}
\tightlist
\item
  STARプロジェクトでの関心：

  \begin{itemize}
  \tightlist
  \item
    \textbf{少人数学級} (処置変数) \(\rightarrow\)
    \textbf{生徒の学業成績} (結果変数)
  \end{itemize}
\end{itemize}
\end{frame}

\begin{frame}[fragile]{処置変数と結果変数の定義}
\protect\hypertarget{ux51e6ux7f6eux5909ux6570ux3068ux7d50ux679cux5909ux6570ux306eux5b9aux7fa9}{}
\begin{itemize}
\tightlist
\item
  \textbf{処置変数 (\(X\))}:

  \begin{itemize}
  \tightlist
  \item
    処置 (\(X_i = 1\)): 処置があるという状態（少人数学級に通った）。
  \item
    統制 (\(X_i = 0\)): 処置がないという状態（標準規模学級に通った）。
  \end{itemize}
\item
  \textbf{結果変数 (\(Y\))}:

  \begin{itemize}
  \tightlist
  \item
    \texttt{reading} (読解力テストの点数)
  \item
    \texttt{math} (算数テストの点数)
  \item
    \texttt{graduated} (高校卒業を示す二値変数: 1=卒業,
    0=卒業しなかった)
  \end{itemize}
\end{itemize}
\end{frame}

\hypertarget{ux500bux5225ux56e0ux679cux52b9ux679c}{%
\section{2.3 個別因果効果}\label{ux500bux5225ux56e0ux679cux52b9ux679c}}

\begin{frame}{潜在的結果と個別因果効果}
\protect\hypertarget{ux6f5cux5728ux7684ux7d50ux679cux3068ux500bux5225ux56e0ux679cux52b9ux679c}{}
\begin{itemize}
\tightlist
\item
  \textbf{潜在的結果}:
  処置を受けた場合と受けなかった場合のそれぞれの結果。

  \begin{itemize}
  \tightlist
  \item
    \(Y_i(X_i = 1)\): 個人\(i\)が処置を受けた場合の潜在的結果。
  \item
    \(Y_i(X_i = 0)\): 個人\(i\)が統制条件下での潜在的結果。
  \end{itemize}
\item
  \textbf{個別因果効果} = \(Y_i(X_i = 1) - Y_i(X_i = 0)\)
\end{itemize}
\end{frame}

\begin{frame}{因果推論の根本問題}
\protect\hypertarget{ux56e0ux679cux63a8ux8ad6ux306eux6839ux672cux554fux984c}{}
\begin{itemize}
\tightlist
\item
  現実の世界では、同じ個人について\textbf{両方の潜在的結果を同時に観察することは決してない}。
\item
  観察できるのは\textbf{現実の結果}のみであり、\textbf{反事実の結果}は観察できない。
\item
  したがって、個人レベルでの因果効果を直接計算することは不可能である。
\end{itemize}
\end{frame}

\hypertarget{ux5e73ux5747ux56e0ux679cux52b9ux679c}{%
\section{2.4 平均因果効果}\label{ux5e73ux5747ux56e0ux679cux52b9ux679c}}

\begin{frame}{平均因果効果（平均処置効果）}
\protect\hypertarget{ux5e73ux5747ux56e0ux679cux52b9ux679cux5e73ux5747ux51e6ux7f6eux52b9ux679c}{}
\begin{itemize}
\item
  個人の効果の代わりに、集団における因果効果の平均に注目する。
\item
  もし全員の両方の潜在的結果を観察できれば：

  \[平均処置効果 \] \[= 個別因果効果の平均 \]
  \[=\overline{Y(X=1)} - \overline{Y(X=0)}\]
\end{itemize}
\end{frame}

\begin{frame}{無作為化実験と平均の差推定量}
\protect\hypertarget{ux7121ux4f5cux70baux5316ux5b9fux9a13ux3068ux5e73ux5747ux306eux5deeux63a8ux5b9aux91cf}{}
\begin{itemize}
\tightlist
\item
  \textbf{無作為化実験}: 処置の割り当てが無作為に行われる研究デザイン。
\item
  無作為化により、処置群と統制群は、処置以外のすべての特徴において\textbf{平均的に互いに同一（比較可能）}になる。
\item
  この場合、現実の結果を用いて平均処置効果を推定できる：

  \begin{itemize}
  \tightlist
  \item
    \textbf{平均の差推定量} =
    \(\overline{Y}_{\text{処置群}} - \overline{Y}_{\text{統制群}}\)
  \end{itemize}
\end{itemize}
\end{frame}

\hypertarget{ux5c11ux4ebaux6570ux5b66ux7d1aux306fux751fux5f92ux306eux6210ux7e3eux3092ux5411ux4e0aux3055ux305bux308bux304b}{%
\section{2.5
少人数学級は生徒の成績を向上させるか？}\label{ux5c11ux4ebaux6570ux5b66ux7d1aux306fux751fux5f92ux306eux6210ux7e3eux3092ux5411ux4e0aux3055ux305bux308bux304b}}

\begin{frame}[fragile]{データの読み込み}
\protect\hypertarget{ux30c7ux30fcux30bfux306eux8aadux307fux8fbcux307f}{}
\begin{Shaded}
\begin{Highlighting}[]
\CommentTok{\# データセットの読み込み}
\NormalTok{star }\OtherTok{\textless{}{-}} \FunctionTok{read.csv}\NormalTok{(}\StringTok{"STAR.csv"}\NormalTok{)}
\end{Highlighting}
\end{Shaded}

\begin{Shaded}
\begin{Highlighting}[]
\CommentTok{\# データの最初の数行を表示}
\FunctionTok{head}\NormalTok{(star)}
\end{Highlighting}
\end{Shaded}

\begin{verbatim}
##   classtype reading math graduated
## 1     small     578  610         1
## 2   regular     612  612         1
## 3   regular     583  606         1
## 4     small     661  648         1
## 5     small     614  636         1
## 6   regular     610  603         0
\end{verbatim}
\end{frame}

\begin{frame}[fragile]{Rの関係演算子と新しい変数の作成}
\protect\hypertarget{rux306eux95a2ux4fc2ux6f14ux7b97ux5b50ux3068ux65b0ux3057ux3044ux5909ux6570ux306eux4f5cux6210}{}
\begin{itemize}
\tightlist
\item
  \texttt{==} を用いて、特定の値と等しいかを評価できる。
\item
  \texttt{ifelse()} 関数を用いて、条件に基づいて新しい変数を作成する。
\end{itemize}

\begin{Shaded}
\begin{Highlighting}[]
\CommentTok{\# 論理テストの例}
\DecValTok{3} \SpecialCharTok{==} \DecValTok{3}
\end{Highlighting}
\end{Shaded}

\begin{verbatim}
## [1] TRUE
\end{verbatim}

\begin{Shaded}
\begin{Highlighting}[]
\DecValTok{3} \SpecialCharTok{==} \DecValTok{4}
\end{Highlighting}
\end{Shaded}

\begin{verbatim}
## [1] FALSE
\end{verbatim}

\begin{Shaded}
\begin{Highlighting}[]
\CommentTok{\# 少人数学級に通ったかを示す新しい二値変数 \textquotesingle{}small\textquotesingle{} を作成}
\NormalTok{star}\SpecialCharTok{$}\NormalTok{small }\OtherTok{\textless{}{-}} \FunctionTok{ifelse}\NormalTok{(star}\SpecialCharTok{$}\NormalTok{classtype }\SpecialCharTok{==} \StringTok{"small"}\NormalTok{, }\DecValTok{1}\NormalTok{, }\DecValTok{0}\NormalTok{)}
\end{Highlighting}
\end{Shaded}
\end{frame}

\begin{frame}[fragile]{新しい変数 `small' を確認する}
\protect\hypertarget{ux65b0ux3057ux3044ux5909ux6570-small-ux3092ux78baux8a8dux3059ux308b}{}
\begin{itemize}
\tightlist
\item
  データの最初の数行を表示
\end{itemize}

\begin{Shaded}
\begin{Highlighting}[]
\FunctionTok{head}\NormalTok{(star)}
\end{Highlighting}
\end{Shaded}

\begin{verbatim}
##   classtype reading math graduated small
## 1     small     578  610         1     1
## 2   regular     612  612         1     0
## 3   regular     583  606         1     0
## 4     small     661  648         1     1
## 5     small     614  636         1     1
## 6   regular     610  603         0     0
\end{verbatim}
\end{frame}

\begin{frame}[fragile]{変数の部分集合化と平均の計算 (1)}
\protect\hypertarget{ux5909ux6570ux306eux90e8ux5206ux96c6ux5408ux5316ux3068ux5e73ux5747ux306eux8a08ux7b97-1}{}
\begin{itemize}
\tightlist
\item
  \texttt{{[}{]}} を用いて、特定の条件を満たす観察を抽出する。
\end{itemize}

\begin{Shaded}
\begin{Highlighting}[]
\CommentTok{\# 読解力の全体平均}
\FunctionTok{mean}\NormalTok{(star}\SpecialCharTok{$}\NormalTok{reading)}
\end{Highlighting}
\end{Shaded}

\begin{verbatim}
## [1] 628.803
\end{verbatim}

\begin{Shaded}
\begin{Highlighting}[]
\CommentTok{\# 処置群（少人数学級）の読解力の平均}
\FunctionTok{mean}\NormalTok{(star}\SpecialCharTok{$}\NormalTok{reading[star}\SpecialCharTok{$}\NormalTok{small }\SpecialCharTok{==} \DecValTok{1}\NormalTok{])}
\end{Highlighting}
\end{Shaded}

\begin{verbatim}
## [1] 632.7026
\end{verbatim}
\end{frame}

\begin{frame}[fragile]{変数の部分集合化と平均の計算 (2)}
\protect\hypertarget{ux5909ux6570ux306eux90e8ux5206ux96c6ux5408ux5316ux3068ux5e73ux5747ux306eux8a08ux7b97-2}{}
\begin{Shaded}
\begin{Highlighting}[]
\CommentTok{\# 統制群（標準規模学級）の読解力の平均}
\FunctionTok{mean}\NormalTok{(star}\SpecialCharTok{$}\NormalTok{reading[star}\SpecialCharTok{$}\NormalTok{small }\SpecialCharTok{==} \DecValTok{0}\NormalTok{])}
\end{Highlighting}
\end{Shaded}

\begin{verbatim}
## [1] 625.492
\end{verbatim}
\end{frame}

\begin{frame}[fragile]{平均の差推定量の計算（読解力と算数）}
\protect\hypertarget{ux5e73ux5747ux306eux5deeux63a8ux5b9aux91cfux306eux8a08ux7b97ux8aadux89e3ux529bux3068ux7b97ux6570}{}
\scriptsize

\begin{Shaded}
\begin{Highlighting}[]
\CommentTok{\# 読解力についての平均の差推定量}
\FunctionTok{mean}\NormalTok{(star}\SpecialCharTok{$}\NormalTok{reading[star}\SpecialCharTok{$}\NormalTok{small }\SpecialCharTok{==} \DecValTok{1}\NormalTok{]) }\SpecialCharTok{{-}} \FunctionTok{mean}\NormalTok{(star}\SpecialCharTok{$}\NormalTok{reading[star}\SpecialCharTok{$}\NormalTok{small }\SpecialCharTok{==} \DecValTok{0}\NormalTok{])}
\end{Highlighting}
\end{Shaded}

\begin{verbatim}
## [1] 7.210547
\end{verbatim}

\begin{Shaded}
\begin{Highlighting}[]
\CommentTok{\# 算数についての平均の差推定量}
\FunctionTok{mean}\NormalTok{(star}\SpecialCharTok{$}\NormalTok{math[star}\SpecialCharTok{$}\NormalTok{small }\SpecialCharTok{==} \DecValTok{1}\NormalTok{]) }\SpecialCharTok{{-}} \FunctionTok{mean}\NormalTok{(star}\SpecialCharTok{$}\NormalTok{math[star}\SpecialCharTok{$}\NormalTok{small }\SpecialCharTok{==} \DecValTok{0}\NormalTok{])}
\end{Highlighting}
\end{Shaded}

\begin{verbatim}
## [1] 5.989905
\end{verbatim}

\normalsize
\end{frame}

\begin{frame}[fragile]{平均の差推定量の計算（卒業率）}
\protect\hypertarget{ux5e73ux5747ux306eux5deeux63a8ux5b9aux91cfux306eux8a08ux7b97ux5352ux696dux7387}{}
\scriptsize

\begin{Shaded}
\begin{Highlighting}[]
\CommentTok{\# 卒業についての平均の差推定量}
\FunctionTok{mean}\NormalTok{(star}\SpecialCharTok{$}\NormalTok{graduated[star}\SpecialCharTok{$}\NormalTok{small }\SpecialCharTok{==} \DecValTok{1}\NormalTok{]) }\SpecialCharTok{{-}} \FunctionTok{mean}\NormalTok{(star}\SpecialCharTok{$}\NormalTok{graduated[star}\SpecialCharTok{$}\NormalTok{small }\SpecialCharTok{==} \DecValTok{0}\NormalTok{])}
\end{Highlighting}
\end{Shaded}

\begin{verbatim}
## [1] 0.007031124
\end{verbatim}

\begin{Shaded}
\begin{Highlighting}[]
\CommentTok{\# 解釈の補助}
\FunctionTok{mean}\NormalTok{(star}\SpecialCharTok{$}\NormalTok{graduated[star}\SpecialCharTok{$}\NormalTok{small }\SpecialCharTok{==} \DecValTok{1}\NormalTok{]) }\CommentTok{\# 処置群の卒業率}
\end{Highlighting}
\end{Shaded}

\begin{verbatim}
## [1] 0.8735043
\end{verbatim}

\begin{Shaded}
\begin{Highlighting}[]
\FunctionTok{mean}\NormalTok{(star}\SpecialCharTok{$}\NormalTok{graduated[star}\SpecialCharTok{$}\NormalTok{small }\SpecialCharTok{==} \DecValTok{0}\NormalTok{]) }\CommentTok{\# 統制群の卒業率}
\end{Highlighting}
\end{Shaded}

\begin{verbatim}
## [1] 0.8664731
\end{verbatim}

\normalsize

\begin{itemize}
\tightlist
\item
  \textbf{解釈}:
  少人数学級に通うことで、高校を卒業する確率が平均で約0.7パーセンテージポイント上昇した。
\end{itemize}
\end{frame}

\hypertarget{ux307eux3068ux3081}{%
\section{2.6 まとめ}\label{ux307eux3068ux3081}}

\begin{frame}[fragile]{第2章のまとめ}
\protect\hypertarget{ux7b2c2ux7ae0ux306eux307eux3068ux3081}{}
\begin{itemize}
\tightlist
\item
  \textbf{因果効果}: 処置変数が結果変数に与える影響。
\item
  \textbf{因果推論の根本問題}: 反事実の結果は観察できない。
\item
  \textbf{無作為化実験}:
  処置群と統制群を比較可能にし、\textbf{平均の差推定量}を用いて平均因果効果を有効に推定できる。
\item
  Rの操作: \texttt{==} 演算子、\texttt{ifelse()}
  による変数の作成、\texttt{{[}{]}} によるデータの部分集合化を学んだ。
\end{itemize}
\end{frame}

\end{document}
